\section{Residue discs at fixed depth}\label{sec:discs}

The proof of Theorem~\ref{thm:discs} uses a relation whose coefficients
may depend on every $p$-adic period occurring in localisation. Its degree
will depend only on $s=|S|$. Ghosh and Ward obtained
$O_n(p/\log p)$ for the roots of a single $\ell_n$ with $n\ge2$
\cite[Theorem~1]{GW}. In the Chabauty--Kim setting,
Betts, Kumpitsch and L\"udtke used a finite polylogarithm of weight
$p-3$ to isolate one Teichm\"uller point
\cite[Proposition~5.15]{BKL}. Here the weights stay fixed and the
relation may mix them. We first bound the number of residues satisfying
any such relation. Write
\[
  \ell_i(X)=\sum_{m=1}^{p-1}\frac{X^m}{m^i}\in\mathbb F_p[X],
  \qquad \boldsymbol\ell=(\ell_1,\ldots,\ell_n).
\]

\begin{lemma}\label{lem:mixed-disc-count}
Fix $n\ge2$ and nonnegative integers $h,d$. For every prime $p$ and every
nonzero $R\in\mathbb F_p[X,y_1,\ldots,y_n]$ with $\deg_X R\le h$
and total $y$-degree at most $d$,
\[
 \#\{a\in\mathbb F_p\setminus\{0,1\}:
                  R(a,\boldsymbol\ell(a))=0\}
       \ll_{n,h,d}p^{1-1/(2n)} .
\]
\end{lemma}

\begin{proof}
We use the auxiliary variables $X,W,y_1,\ldots,y_n$ in the manner of
Ghosh and Ward \cite[\S3]{GW}. Put
\[
 \delta=X(1-X)\partial_X+(X-W)\partial_{y_1}
             +(1-X)\sum_{i=2}^n y_{i-1}\partial_{y_i},
 \qquad \delta W=0.
\]
Upon substituting $W=X^p$ and $y_i=\ell_i(X)$ this derivation induces
$D=X(1-X)\,d/dX$. Indeed $X(1-X)\ell_1'=X-X^p$ by summing a geometric
progression, and $X\ell_i'=\ell_{i-1}$ for $i\ge2$.

Set $N_t=\binom{t+n}{n}$, choose an integer $b>3nd+h+1$, and put
$A=r=t^2$, $B=bt$. Seek $\Phi$ in the space of polynomials of
degrees at most $A,B,t$ in $X,W,\boldsymbol y$, respectively, taking
total degree in $\boldsymbol y$. For $0\le j<r$ impose the linear conditions
\[
             R\mid (\delta^j\Phi)(X,X,\boldsymbol y).
\]
The right-hand polynomials have $X$-degree at most $E=A+B+r$ and
$y$-degree at most $t$, since $\delta$ raises total $(X,W)$-degree by
at most one and preserves the $y$-degree bound. In this space, multiples
of $R$ include the injective image under multiplication by $R$ of the
box with bounds $E-h,t-d$. For $t$ sufficiently large the codimension is
therefore at most
\[
 (E+1)(N_t-N_{t-d})+hN_{t-d}
       \le (3ndt+h)N_t .
\]
Here $E+1\le 3t^2$ and telescoping the last $d$ binomial increments gives
$N_t-N_{t-d}\le d\binom{t+n-1}{n-1}\le dnN_t/(t+1)$.
The total codimension is smaller than $(A+1)(B+1)N_t$, so a nonzero
$\Phi$ exists.

Choose $t=\lfloor c_n p^{1/(2n)}\rfloor$ with $c_n>0$ sufficiently small.
For $A=t^2$, the bound $N_t A+\binom{N_t}{2}=O_n(t^{2n})$ uses $n\ge2$.
Thus for large $p$ the bound is strictly below $p$, and $t,r<p$.
Corollary~\ref{cor:box} now shows that for every nonzero
coefficient $f(X,\boldsymbol y)$ of $\Phi$ as a polynomial in $W$,
\[
                       X^p\nmid f(X,\boldsymbol\ell(X)).
\]
In particular, taking the first nonzero coefficient, with index $i$,
the corresponding summand of
$F(X)=\Phi(X,X^p,\boldsymbol\ell(X))$ has order strictly below $(i+1)p$
at zero, and all subsequent summands have order at least $(i+1)p$.
Hence $F\ne0$ and $\deg F\le A+pB+(p-1)t$.

At a residue $a$ counted by the lemma, $a^p=a$, so the divisibility
conditions give $(D^jF)(a)=0$ for all $j<r$. If $F$ had order $m<r$
there, with leading coefficient $u$ in $X-a$, we would have
$(D^mF)(a)=m![a(1-a)]^mu\ne0$. Thus each residue contributes at
least $r$ to the degree of $F$, giving an upper bound
$1+(b+1)p/t$. The finitely many smaller primes can be absorbed into the
implied constant.
\end{proof}

\begin{proof}[Proof of Theorem~\ref{thm:discs}]
Use the standard integral model
$X=\mathbb P^1_{\mathbb Z}\setminus\{0,1,\infty\}$ and its usual
tangential base point at zero directed towards one for polylogarithmic
coordinates. Take $p\notin S$ and $s\ge1$. We may suppose $p>n+1$,
where $n=2s$ for the refined locus and $n=4s$ for the unrefined locus.
For smaller $p$ the trivial bound $p-2$ suffices, with a constant
depending only on $s$. Passing from the full depth-$n$ fundamental
group to its polylogarithmic quotient carries the global and local Kummer
conditions into the corresponding quotient conditions. It suffices to bound
the reductions of the resulting polylogarithmic locus.

The localisation formula of Betts, Kumpitsch and L\"udtke
\cite[Propositions~5.3 and~5.9, Theorem~5.6]{BKL}
has global coordinates $x_\ell,y_\ell$ for $\ell\in S$, of weight one,
and $z_j$ for odd $3\le j<n$, of weight $j$ (here $n$ is even).
In the refined case $2\notin S$ forces emptiness by the local
condition $X(\mathbb Z_2)=\varnothing$. If $2\in S$, there are $3^s$
refinement conditions, choosing for each $\ell$ one of the hyperplanes
$x_\ell=0$, $y_\ell=0$, $x_\ell+y_\ell=0$
\cite[Proposition~5.11]{BKL}. On each resulting linear subspace there
are $s+(s-1)=n-1$ free coordinates. In the unrefined case the whole
source has $2s+(2s-1)=n-1$ coordinates. Denote either list by $v$.
In these coordinates the local functions
\[
 L=p^{-1}\log_p,\qquad Y_m=p^{-m}\operatorname{Li}_m\quad(1\le m\le n)
\]
pull back to polynomials $L(v),Y_m(v)$ over $K=\mathbb Q_p$ of ordinary
degrees at most $1,m$, respectively. In fact each term in the
localisation formula of weight $m$ is a product of $m$ weight-one
coordinates, or of a coordinate $z_j$ and $m-j$ weight-one coordinates;
restricting to the chosen linear subspace preserves the degree bound.

Adjoin a free variable $x$. Define $W_0=-x$, $W_m=(x-1)Y_m(v)$ for $m>0$,
and form the map from $\mathbb A_K^n$ to $\mathbb A_K^{n+1}$ with coordinates
\[
 x,\qquad u_m=\sum_{j=0}^{m}W_{m-j}\frac{(-L(v))^j}{j!}
           \quad (1\le m\le n).
\]
All coordinates have degree at most $D=n+1$. Set $H=((n+1)D+1)^n$.
Polynomials of individual variable degrees at most $H$ in the target
have dimension $(H+1)^{n+1}$. After substitution their individual
variable degrees in the source are at most $(n+1)DH$. The source
space has dimension at most
\[
 ((n+1)DH+1)^n\le H(H+1)^n<(H+1)^{n+1}.
\]
Consequently a nonzero polynomial $P$ of individual degrees at most $H$
vanishes identically after substitution. Divide by a coefficient of
maximal $p$-adic absolute value to make $P$ integral with nonzero
reduction $R$. This normalization uses only coefficients of $P$ and
requires no integrality of $v$ or of the periods. Moreover,
$P(x,\boldsymbol u)$, viewed as a polynomial in $x$ and the local
coordinates $L,\boldsymbol Y$, vanishes on the product of the $x$-line
with the Zariski closure of the localisation image. This follows by
taking its coefficients in $x$, each of which pulls back to zero.

We compute $\boldsymbol u$ modulo $p$ at a local point $z$ reducing
to $a\ne0,1$. Let $\zeta$ be the Teichm\"uller lift of $a$, and write
$z=\zeta(1+pw)$ with $w\in\mathbb Z_p$. Then
$L(z)\equiv w\pmod p$. At $\zeta$, the $p$-depleted polylogarithm
$\operatorname{Li}_m(\zeta)-p^{-m}\operatorname{Li}_m(\zeta^p)$
equals $(p^m-1)Y_m(\zeta)$. Its reduction is
$\ell_m(a)/(1-a)$ by Besser \cite[Proposition~2.1]{Besser}. Hence
$Y_m(\zeta)$ is integral with reduction $\ell_m(a)/(a-1)$.
Set $Y_0(\zeta)=\zeta/(1-\zeta)$. Besser's translation formula
\cite[Proposition~2.3]{Besser} gives integral $Y_m(z)$ and
\[
 Y_m(z)\equiv\sum_{j=0}^m Y_{m-j}(\zeta)w^j/j!\pmod p .
\]
Take any integral $x$ lifting $a$. Since $(a-1)Y_0(\zeta)\equiv-a$,
the reductions of $W_0,W_1,\ldots,W_n$ are the coefficients through
degree $n$ of
\[
 \left(-a+\sum_{m=1}^n\ell_m(a)T^m\right)
                 \sum_{j=0}^n\frac{\bar w^j T^j}{j!}
                \quad\text{in }\mathbb F_p[T]/(T^{n+1}).
\]
The sum defining $u_m$ multiplies this expression by the truncated
exponential of $-\bar w T$. The two exponentials multiply to $1$
in this ring, by the binomial identity and $n<p$. We obtain
$u_m\equiv\ell_m(a)\pmod p$.

If $z$ lies in the locus under consideration, its local coordinates lie
in the localisation image or its closure, so
$R(a,\boldsymbol\ell(a))=0$. Apply Lemma~\ref{lem:mixed-disc-count} with
$h=H$ and $d=nH$. Each image contributes $O_n(p^{1-1/(2n)})$ discs.
Summing over at most $3^s$ images in the refined case, and taking the
single image in the unrefined case, gives the asserted exponents
$1-1/(4s)$ and $1-1/(8s)$.
\end{proof}

\paragraph{An affine genus-two pullback.}
Let $C$ be the smooth projective model of $y^2=x^5-x$ and let $U$ be
the inverse image of $X$ under the hyperelliptic projection. The branch
divisor consists of $0,\infty$ and the four fourth roots of unity,
which are pairwise distinct in odd characteristic. Thus the double
cover has genus two by Riemann--Hurwitz and good reduction away from
$2$. Removing the three ramification points above $0,1,\infty$ gives
a finite degree-two morphism $U\to X$ over $\mathbb Z[1/2]$.
Use the compatible integral base points $(-1,0)$ and $-1$. Base-point
change on $X$ twists path torsors by the path torsor between the two
base points and identifies the corresponding local and global Kummer
conditions; thus the disc bound computed in tangential coordinates
applies with base point $-1$ as well.

For $S=\{2\}$, functoriality sends the full refined depth-two locus of
$U$ into that of $X$: the homomorphism of unipotent fundamental groups
preserves the lower central series, and its induced map on torsors
commutes with localisation and the local Kummer maps. The image of
each local Kummer closure lies in the corresponding closure for $X$,
by continuity in the Zariski topology. At an odd prime $p$, each
residue $a$ on the target has at most two lifts, since its fibre
satisfies $y^2=a^5-a$. Theorem~\ref{thm:discs} gives
$O(p^{3/4})$ occupied discs for the full refined depth-two locus
of $U$. This conclusion concerns the open curve with the specified
three punctures.
